% !TeX root = ../main.tex

\chapter{绪论}
\label{chap:intro}

在线图推理的核心矛盾在于，多层消息传递虽然能够提升节点表示能力，但多跳邻域的持续扩张也会显著推高推理代价。针对这一矛盾，稀疏分解图神经网络（Sparse Decomposition of Graph Neural Networks，SDGNN）以单次稀疏分解替代逐层消息传递，为高效在线推理提供了可行思路。然而，其全局统一的稀疏预算机制仍难以适应真实图中普遍存在的节点拓扑异质性。围绕这一局限，本文提出多尺度自适应稀疏图神经网络（Multi-Scale Adaptive Sparse Graph Neural Network，MSAS-GNN）\label{idx:name:msasgnn}。本章将从研究背景与意义、国内外研究现状、研究内容与创新点，以及技术路线与论文组织结构四个方面，说明本文的研究动机、问题来源与整体安排。与方法相关的核心符号、形式化定义和理论分析将在第~3--5~章中进一步展开。

\begin{figure}
  \centering
  \includegraphics[width=1\linewidth]{1-1.pdf}
  % \setlength{\fboxsep}{1.2em}
  % \fbox{\parbox{0.88\textwidth}{\small\centering
  %   \textbf{【图1.1 占位符，定稿前须替换为正式图】}\\[0.6em]
  %   图题：SDGNN全局统一稀疏预算局限与MSAS-GNN谱驱动自适应改进示意图\\[0.4em]
  %   表达内容：图采用左右对照结构。左侧以统一稀疏预算$\lambda_1$同时作用于高度数枢纽节点与低度数外围节点，突出“一刀切”分配带来的两类失配：前者易丢失关键邻居，后者保留较多冗余连接；右侧展示MSAS-GNN依据节点谱域特征进行差异化预算分配，在同配与异配两类图上分别保留更匹配的邻域信息。图中仅保留概念层面的模块框、节点示意与箭头关系，不出现后续章节的具体公式与算法编号。
  % }}
  \caption{SDGNN全局统一稀疏预算局限与MSAS-GNN谱驱动自适应改进示意图}
  \label{fig:ch1-limitation-overview}
\end{figure}


%% ============================================================
\section{研究背景与意义}
\label{sec:ch1-background}

图结构数据以节点和边刻画实体及其关联关系，广泛存在于社交网络、推荐系统、知识图谱、生物分子网络与交通网络等场景中。与图像、文本等规则结构数据不同，图数据通常具有邻接关系不规则、局部结构差异显著、节点邻域规模不一等特点，传统面向欧式数据的深度学习模型难以直接利用这类拓扑信息\cite{JSJX202201003}。围绕这一问题发展起来的图神经网络（Graph Neural Network，GNN）通过在图拓扑上进行消息传递与邻域聚合，逐层融合节点自身特征与多跳邻域信息，在节点分类、链接预测、图分类等任务上表现出较强的建模能力，因而成为图表示学习的重要方法。

尽管GNN在多类任务上取得了良好效果，但其推理阶段仍面临明显的效率瓶颈。对于$L$层消息传递模型，目标节点的有效感受野会随层数扩大，在平均度为$\bar{d}$的图上，其关联节点规模通常可近似写为$O(\bar{d}^{L})$。在这一过程中，邻居检索、特征读取与多轮聚合带来的计算和内存访问开销会迅速增加，使得GNN在大规模图和低时延场景下的部署受到限制。尤其在推荐系统、实时风控等在线服务场景中，模型不仅需要保持表示能力，还必须满足毫秒级响应要求，因此如何在尽量保留图结构信息的前提下降低推理代价，成为图神经网络走向实际应用时必须面对的问题。

针对上述问题，Hu等提出稀疏分解图神经网络（Sparse Decomposition of Graph Neural Networks，SDGNN\label{idx:name:sdgnn}）\cite{Hu2024SDGNN}，将教师GNN表示的逼近重写为单次稀疏分解\label{idx:theme:sparsedecomp}。在该框架中，推理阶段仅需执行
\begin{equation}
  \hat{H}=\Theta^\top\tilde{\Phi}.
\end{equation}
其中，$\tilde{\Phi}\in\mathbb{R}^{n\times d}$表示经过特征变换并作行归一化后的节点表示矩阵，其第$i$行对应节点$i$的变换特征；$\Theta\in\mathbb{R}^{n\times n}$为稀疏权重矩阵，第$i$列$\theta_i$刻画节点$i$对其候选邻居特征的线性组合权重；$\hat{H}\in\mathbb{R}^{n\times d}$为稀疏分解得到的节点表示矩阵。由于每个节点最终仅保留少量非零权重，若记全图平均非零权重数为$\bar{k}$，则该推理过程的主项复杂度可降至$O(n\bar{k}d)$。与逐层消息传递相比，这种做法避免了推理时反复展开多跳邻域，为GNN的高效在线推理提供了新的实现路径。与此同时，SDGNN通过LARS求解器在全局统一正则系数$\lambda_{\mathrm{reg}}$下构造Lasso子问题，并在训练结束后固化稀疏权重矩阵$\Theta$，从而兼顾了近似精度与推理效率。

然而，以同一$\lambda_{\mathrm{reg}}$约束所有节点的稀疏预算\label{idx:theme:sparsebud}，实际上隐含了“不同节点适合相同稀疏压力”的前提。真实图中节点的拓扑角色通常存在显著差异\label{idx:theme:tophete}：枢纽节点往往连接更广，承担更多跨区域信息传递功能，若稀疏化过强，容易丢失关键邻居；外围节点的局部冗余则相对更多，若稀疏化不足，又会削弱压缩收益。在结构更复杂、局部频谱参与方式差异更明显的图中，这种全局统一参数更容易与节点的真实信息需求发生失配，如图~\ref{fig:ch1-limitation-overview}~所示。由此可见，如何在保持$O(n\bar{k}d)$推理口径不变的前提下，将节点间的结构差异引入稀疏预算分配过程，是现有稀疏分解路线进一步发展的关键问题。

基于这一观察，本文尝试从图谱理论出发，建立面向节点差异的自适应稀疏预算分配框架。具体而言，本文希望以可计算的图结构统计量为依据，将节点的谱域参与特征、局部结构复杂性与多跳传播差异转化为可解释的稀疏参数设计规则，使稀疏化过程不再完全依赖全局统一预算，而能够对不同结构角色的节点采取更有针对性的分配方式。该研究一方面有助于为稀疏推理中的节点级预算设计提供更清晰的理论依据，另一方面也为大规模静态图上的高效部署提供更具适配性的实现思路。本文方法的具体效果将在第~6~章中通过实验进行验证。

\section{国内外研究现状}
\label{sec:ch1-related}

围绕本文关注的“高效在线图推理”问题，现有研究大致可以归纳为三条彼此相关的文献脉络：其一是图神经网络推理效率及其瓶颈分析，其二是面向效率提升的高效与稀疏图神经网络方法，其三是复杂图结构学习中的谱视角方法。以下仅结合本文研究问题作概要梳理，为第~\ref{sec:ch1-contribution}~节的问题界定和第~2~章的系统综述提供铺垫。

\subsection{图神经网络在线推理效率相关研究}
\label{subsec:ch1-gnn-efficiency}

Gilmer等\cite{Gilmer2017MPNN}\label{idx:name:gilmer}将多种图学习方法统一为消息传递神经网络（MPNN）框架，之后的GCN\cite{Kipf2017GCN}\label{idx:name:gcn}、GAT\cite{Velickovic2018GAT}\label{idx:name:gat}和GIN\cite{Xu2019GIN}\label{idx:name:gin}等模型，分别从谱滤波、注意力加权和表达能力等角度扩展了这一范式。此类方法的共同特点是通过多层邻域聚合逐步扩大节点的感受野，从而提升表示能力，但其代价是推理计算开销会随层数、边数和隐藏维度同步增长。以GCN为例，在全图一次前向的口径下，其主项复杂度可写为$O(L\cdot m\cdot d)$。

对于离线批处理场景，上述复杂度仍具有可接受性；但在推荐系统、实时风控等强调低时延响应的在线服务中，目标往往是对少量查询节点快速完成推理，此时多跳邻居收集与内存访问会成为主要瓶颈。与此同时，随着传播层数增加，模型还可能出现过平滑\label{idx:theme:oversmooth}\cite{Li2018OverSmoothing}和过压缩\label{idx:theme:oversquash}\cite{Alon2021OverSquashing}等现象，使得简单地依靠加深网络来换取更大感受野并不可取。这些问题共同说明，高效在线推理并不仅仅是“减少计算量”那么简单，还需要重新审视GNN表示形成过程中的信息保留方式，从而为稀疏分解等替代路径提供了问题背景。

\subsection{高效与稀疏图神经网络方法}
\label{subsec:ch1-sparse}

针对GNN推理中的感受野扩张问题，现有方法大致沿邻居采样、预计算线性化和稀疏分解三条路线展开。邻居采样方法，如GraphSAGE\cite{Hamilton2017GraphSAGE}\label{idx:name:graphsage}、FastGCN\cite{Chen2018FastGCN}\label{idx:name:fastgcn}和GraphSAINT\cite{Zeng2020GraphSAINT}\label{idx:name:graphsaint}，通过限制每层参与聚合的邻居规模来降低训练成本，但在推理阶段仍往往需要访问较大范围的邻域，且采样策略通常较为固定，难以针对不同节点的结构差异进行精细调整。

预计算线性化方法，如SGC\cite{Wu2019SGC}\label{idx:name:sgc}、APPNP\cite{Gasteiger2019APPNP}\label{idx:name:appnp}和SIGN\cite{Frasca2020SIGN}，则通过提前完成部分传播计算，将推理阶段尽量简化为轻量级线性变换，从而显著降低在线代价。这一路线在效率上具有明显优势，但传播模式通常在预处理阶段已经固定，因而缺少节点级自适应能力。

与前两类方法相比，稀疏分解路线更直接地面向“用少量关键邻居逼近教师表示”这一目标。SDGNN\cite{Hu2024SDGNN}在这一方向上给出了代表性方案，其将推理主项复杂度压缩到$O(n\bar{k}d)$（$\bar{k}\ll\bar{d}^{L}$），成为本文的直接比较基线。需要指出的是，SDGNN的优势主要来自推理阶段对多层传播过程的显式替代，而其不足则集中体现在预算控制机制上：全局统一正则系数$\lambda_{\mathrm{reg}}$对所有节点施加相同稀疏压力，难以反映节点在拓扑角色、局部结构和频域参与方式上的差异。这一局限正是本文在第~3--5~章中进一步展开自适应稀疏机制设计的直接出发点。

\subsection{复杂图结构学习与谱视角方法}
\label{subsec:ch1-heterophily}

除推理效率问题外，图结构本身的复杂性也会影响GNN的有效建模方式。已有研究表明，不同结构类型的图在判别信息所处的有效频段上可能存在明显差异：结构较平滑的引文网络通常更依赖低频聚合，而网页协作图等复杂结构数据则往往需要保留更多中高频成分\cite{Pei2020GeomGCN}。这一现象说明，统一的低通传播假设并不总是适合所有图数据。

围绕复杂图结构建模，H2GCN\cite{Zhu2020H2GCN}\label{idx:name:h2gcn}、GPR-GNN\cite{Chien2021GPRGNN}\label{idx:name:gprgnn}、FAGCN\cite{Bo2021FAGCN}\label{idx:name:fagcn}和LINKX\cite{Lim2021LINKX}\label{idx:name:linkx}等方法，分别从高阶邻域分离、自适应传播权重、高低频成分联合建模以及特征变换与传播解耦等角度，提高了模型对复杂图结构的适应能力。这些方法表明，谱视角能够为理解图结构差异及其传播行为提供有价值的分析工具。

不过，从本文关注的稀疏推理问题看，上述方法大多仍采用全图共享的频率参数或统一传播机制，尚未进一步落实到“节点级频率差异如何转化为节点级稀疏预算”这一层面。谱间隙$\lambda_{\text{gap}}$\label{idx:theme:specgap}能够刻画图的全局连通性，但若要进一步区分不同节点在局部结构中的频率参与特征，还需要引入更细粒度的局部谱统计量。基于这一考虑，本文保留“频率接口”这一设计位置，但在正文主线中采用等权实例化，以保证方法不依赖标签统计量并保持任务通用性；更精细的频率感知设计则作为后续扩展方向保留。

\section{研究内容与创新点}
\label{sec:ch1-contribution}

上述三条文献脉络表明，现有高效图推理方法虽然在不同程度上缓解了GNN的部署开销，但在节点级差异化稀疏预算设计方面仍存在明显不足。以SDGNN为代表的稀疏分解方法主要依赖全局统一参数控制节点稀疏强度，尚未将节点拓扑差异及其频域影响系统地纳入预算设计过程；复杂图结构学习中的谱视角方法则更多停留在全图共享传播参数层面，尚未与节点级稀疏推理机制形成直接结合。这一研究空白构成了本文的核心问题背景。

围绕这一问题，本文关注如下研究目标：在保持静态推理主项复杂度$O(n\bar{k}d)$不变的前提下，如何利用可计算的图结构统计量构建节点级差异化稀疏预算分配方案，使不同结构角色的节点能够获得更匹配的信息保留强度，并使这一设计具有可解释、可分析和可验证的依据。围绕该目标，本文将第~3--5~章组织为一条连续的方法链，依次完成问题形式化与指标提取、三维自适应机制设计，以及算法整合与理论分析。

在方法设计层面\label{idx:innov:3dim}，本文的主要工作可概括为以下三个方面。第一，第~3~章构建了由谱间隙、节点谱能量、归一化局部图熵和节点中心性组成的图复杂度指标体系，为后续自适应参数设计提供统一的可计算输入。第二，第~4~章从频率、节点和跳距三个维度构建自适应稀疏机制：在频率维保留统一频率接口，并采用等权实例化作为正文主线默认配置；在节点维设计节点级正则系数，使不同节点能够对应不同稀疏强度；在跳距维结合多跳传播中的信息衰减现象，设计分层预算分配策略。第三，第~5~章将上述三维参数整合到LARS求解框架中，给出完整算法流程，并进一步分析误差与复杂度，使方法实现与理论讨论形成闭环。在这一设计下，推理阶段仍保持为单次稀疏矩阵乘形式，其主项复杂度维持在线性量级$O(n\bar{k}d)$。

在实验验证层面，本文在六个静态节点分类基准数据集上对所提方法进行了系统评估。实验结果表明，MSAS-GNN相较基线SDGNN在引文网络、网页图和大规模图场景下均表现出稳定改进，节点级自适应正则系数与分层跳距预算均对性能提升具有实际作用。与此同时，方法在ogbn-arxiv等数据集上仍保持较高推理效率，相对消息传递基线取得了明显加速，相关差异也通过Wilcoxon符号秩检验进行了统计验证。更具体的数据、设置与分析将在第~6~章中统一给出。

\section{技术路线与论文组织结构}
\label{sec:ch1-structure}

如图~\ref{fig:ch1-tech-route}~所示，全文共分为七章，可归纳为四个相互衔接的功能层次。第~1--2~章用于建立研究背景、问题动机和理论基础；第~3--5~章构成本文的核心方法链，按照“问题形式化与指标提取$\rightarrow$三维自适应机制设计$\rightarrow$算法整合与理论分析”的顺序展开，并统一以$O(n\bar{k}d)$推理口径为约束；第~6~章对方法进行实验评估；第~7~章对全文工作进行总结，并讨论研究局限与未来方向。

具体而言，第~1~章阐述研究背景、研究现状、研究问题与全文结构；第~2~章围绕高效图推理与谱视角方法梳理相关工作，明确现有研究空白；第~3~章对问题进行形式化描述，并建立后续章节所需的图复杂度指标体系；第~4~章构建频率维、节点维和跳距维三类自适应参数；第~5~章完成算法整合、误差分析与复杂度分析；第~6~章在多类静态节点分类数据集上开展主实验、消融实验和效率评估；第~7~章总结全文工作，并给出后续研究展望。上述安排形成了从问题提出、方法设计到实验验证的完整论证链条。

\section{本章小结}
\label{sec:ch1-summary}

本章的任务在于交代全文研究的出发点，而不是提前展开后续方法细节。围绕在线图推理中的效率与表示能力之间的矛盾，本章首先以SDGNN为切入点，指出全局统一稀疏预算难以适应节点拓扑差异这一核心局限；随后从推理效率、高效稀疏化以及复杂图结构学习三个方面梳理了与本文直接相关的研究进展，并据此归纳出本文关注的关键空白，即节点级自适应稀疏预算设计仍缺少统一的谱论依据和与稀疏推理机制相协调的分析框架。在此基础上，本文将研究目标限定为：在保持$O(n\bar{k}d)$静态推理主项复杂度口径不变的前提下，构建谱驱动的三维自适应稀疏参数体系。下一章将在此基础上进一步系统整理相关理论基础和代表性方法，为后续方法章节展开提供文献与概念支撑。

\begin{figure}[htbp]
  \centering
    \includegraphics[width=1\linewidth]{1-2.pdf}
  % \fbox{\parbox{0.90\linewidth}{
  % 图1.2占位内容：全文技术路线与章节接口图。上层为第~1--2~章的研究背景与理论基础，
  % 中层为第~3--5~章三段式方法链（问题形式化与指标体系$\rightarrow$三维自适应机制设计$\rightarrow$算法整合、理论保证与复杂度分析），
  % 下层为第~6~章实验评估与第~7~章结论展望。箭头仅标示章间输入输出关系，不出现具体公式与算法编号。}}
  \caption{本文技术路线与章节接口示意图}
  \label{fig:ch1-tech-route}
\end{figure}